\subsection{CONNECTED COMMUTATIVE REAL LIE GROUPS}

Let \(G\) be a connected commutative (real) Lie group. The exponential map

\[
\exp_ {\mathrm{G}}: \mathrm{L} (\mathrm{G}) \to \mathrm{G}
\]

is a morphism of Lie groups, surjective with discrete kernel (Chap. III, §6, no. 4, Prop. 11), hence the fact that \(\mathrm{L}(\mathrm{G})\) is a connected covering of \(\mathrm{G}\) .

\begin{enumerate}
\def\labelenumi{\alph{enumi})}
\item
  The following conditions are equivalent: G is simply-connected, \(\exp_{G}\) is an isomorphism, G is isomorphic to \(R^{n}\) ( \(n = \dim G\) ). In this case, transporting the vector space structure of L(G) to G by the isomorphism \(\exp_{G}\) gives a vector space structure on G, which is the only one compatible with the topological group structure of G. Simply-connected commutative Lie groups are called vector (Lie) groups; unless stated otherwise, they are always given the R-vector space structure defined above.
\item
  Denote by \(\Gamma(\mathrm{G})\) the kernel of \(\exp_{G}\) . By General Topology, Chap. VII, §1, no. 1, Th. 1, the group G is compact if and only if \(\Gamma(\mathrm{G})\) is a lattice in \(\mathrm{L}(\mathrm{G})\) , in other words (loc. cit.) if the rank of the free Z-module \(\Gamma(\mathrm{G})\) is equal to the dimension of G. Conversely, if L is a finite dimensional R-vector space and \(\Gamma\) a lattice in L, the quotient topological group \(L/\Gamma\) is a compact connected commutative Lie group.
\end{enumerate}

The compact connected commutative Lie groups are called real tori, or (in this chapter) tori.

\begin{enumerate}
\def\labelenumi{\alph{enumi})}
\setcounter{enumi}{2}
\tightlist
\item
  In the general case, let E be the vector subspace of L(G) generated by \(\Gamma(\mathrm{G})\) , and let V be a complementary subspace. Then G is the direct product of its Lie subgroups \(\exp(\mathrm{E})\) and \(\exp(\mathrm{V})\) ; the first is a torus, the second is vector. Finally, every compact subgroup of G is contained in \(\exp(\mathrm{E})\) (since its projection onto \(\exp(V)\) is necessarily reduced to the identity element); thus, the subgroup \(\exp(E)\) is the unique maximal compact subgroup of G.
\end{enumerate}

For example, take \(\mathrm{G} = \mathrm{C}^*\) ; identify \(\mathrm{L(G)}\) with \(\mathbf{C}\) so that the exponential map of \(\mathrm{G}\) is \(x \mapsto e^x\) . Then \(\Gamma(\mathrm{G}) = 2\pi i\mathbf{Z}, \mathrm{E} = i\mathbf{R}\) , and so \(\exp(\mathrm{E}) = \mathbf{U}\) ; if we take \(\mathrm{V} = \mathbf{R}\) , then \(\exp(\mathrm{V}) = \mathbf{R}_+^*\) and we recover the isomorphism \(\mathbf{C}^* \to \mathbf{U} \times \mathbf{R}_+^*\) constructed in General Topology, Chap. VIII, §1, no. 3, Prop. 1.

\(d\) ) Note finally that \(\exp_{\mathrm{G}}:\mathrm{L}(\mathrm{G})\to \mathrm{G}\) is a universal covering of \(\mathbf{G}\) , hence \(\Gamma (\mathbf{G})\) can be identified naturally with the fundamental group of \(\mathbf{G}\) .

